%% %%%%%%%%%%%%%%%%%%%%%%%%%%%%%%%%%%%%%%%%%%%%%%%%%
%% Template for working papers
%% Tommaso Di Francesco
%% University of Amsterdam
%% %%%%%%%%%%%%%%%%%%%%%%%%%%%%%%%%%%%%%%%%%%%%%%%%%
%% Version 1.0 // February 2024
%% 
\documentclass[11pt]{article}
\usepackage{style}


%% ===============================================
%% Setting the line spacing (3 options: only pick one)
% \doublespacing
% \singlespacing
\onehalfspacing
%% ===============================================

\setlength{\droptitle}{-5em} %% Don't touch

% %%%%%%%%%%%%%%%%%%%%%%%%%%%%%%%%%%%%%%%%%%%%%%%%%%%%%%%%%%
% SET THE TITLE
% %%%%%%%%%%%%%%%%%%%%%%%%%%%%%%%%%%%%%%%%%%%%%%%%%%%%%%%%%%
% TITLE:
\title{Lasso Trading
}
\author{}
% DATE:
\date{\today}

% %%%%%%%%%%%%%%%%%%%%%%%%%%%%%%%%%%%%%%%%%%%%%%%%%%%%%%%%%%
% %%%%%%%%%%%%%%%%%%%%%%%%%%%%%%%%%%%%%%%%%%%%%%%%%%%%%%%%%%
\begin{document}
% %%%%%%%%%%%%%%%%%%%%%%%%%%%%%%%%%%%%%%%%%%%%%%%%%%%%%%%%%%
% %%%%%%%%%%%%%%%%%%%%%%%%%%%%%%%%%%%%%%%%%%%%%%%%%%%%%%%%%%
% ABSTRACT
% %%%%%%%%%%%%%%%%%%%%%%%%%%%%%%%%%%%%%%%%%%%%%%%%%%%%%%%%%%
% %%%%%%%%%%%%%%%%%%%%%%%%%%%%%%%%%%%%%%%%%%%%%%%%%%%%%%%%%%
{\setstretch{.8}
\maketitle
\centering
% %%%%%%%%%%%%%%%%%%


}


% --------------------
\section{Draft}
% --------------------

Consider a simple asset pricing model in which the equilibrium price of a stochastic asset evolves accrodingly to the following law of Motion
\begin{equation}
    p_t = p^{e}_{t+1} + u_t,
\end{equation}

where $p_t$ is the equilibrium price at time $t$, $p^{e}_{t+1}$ is the expected price at time $t+1$ formed at time $t$ \textit{before} the realization of the equilibrium price and $u_t \sim N(0, \sigma^2_u)$.
We assume that agents are boundedly rational and are not aware of the true data generating process.
Instead they rely on econometric techniques to form their expectations.
Specifically we assume that agents have a linear and stationary perceived law of motion (PLM) for the asset return 
\begin{equation}
r_{t+1} = \beta \theta_t + \varepsilon_t,
\end{equation}
where $\theta_t$ is a stochastic signal that agents observe and $\varepsilon_t \sim N(0, \sigma^2_\varepsilon)$.

In order to estimate the parameter agent use Recursive Ordinary Least square of the form 

\begin{equation}
\hat{\beta_t} = \frac{\sum^{t-1}_{s=0} r_{s} \theta_{s-1}}{\sum^{t-1}_{s=0}(\theta_{s-1})^2},
\end{equation}

in which the time indexes are consistent with the timeline of events which implies that at time $t$ the last observable return that can be used for estimation is $r_{t-1}$.
The Actual Law of Motion (ALM) for the asset price obtains by substituting the PLM into the equilibrium condition

\begin{equation}
p_t = p^{e}_{t+1} + u_t = (1 + \hat{\beta_t} \theta_t) p^e_{t},
\end{equation}

where again $p^e_{t}$ stresses that the current price is not available to agents when making their forecats.
Consisten with their model

\begin{equation}
p^e_t = (1 +\hat{\beta_t} \theta_{t+1}) p_{t-1},
\end{equation}
where the use of $\hat{\beta_t}$ is justified by the belief that there is a constant homogenous $\beta$ and so it make sense for agents to use the 'smoothed' estimate obtained by considering also the return at time $t-1$ in place of $\hat{\beta_{t-1}}$.
This implies an ALM for the current price

\begin{equation}
p_t = (1 + \hat{\beta_t} \theta_t)(1 + \hat{\beta_t} \theta_{t-1}) p_{t-1} + u_t,
\end{equation}

and it follows that future returns evolve according to 

\begin{equation}
r_{t+1} = \hat{\beta_{t+1}} \theta_{t+1} + \hat{\beta_{t+1}} \theta_t + (\hat{\beta_{t+1}})^2 \theta_{t+1} \theta_t + \frac{u_t}{p_{t-1}}.
\end{equation}

\subsection{Existence of self-fulfilling equilibria}

A self-fulfilling equilibrium requires consistency between the regression coefficient that agents would estimate from the ALM and the actual $\beta$ implied by their forecast. 
Formally, let the regression coefficient that would be estimated from the ALM if agents used a constant $\beta$ be denoted by $\hat{\beta}$, we require

\begin{equation}
\hat{\beta_t} =  \beta(\hat{\beta_t}) = \frac{\operatorname{Cov}(r_{t+1},\theta_t)}{\operatorname{Var}(\theta_t)}
\end{equation}

where the notation stresses that the population parameter $\beta$ is a function of the estimated parameter $\hat{\beta_t}$.
A self-fulfilling equilibrium is a fixed point of this function, i.e. a solution to the equation $\hat{\beta} = \beta(\hat{\beta})$.

Now
\begin{equation}
\operatorname{Cov}\big(r_{t+1},\theta_t\big)
= \hat{\beta}\,\operatorname{Cov}(\theta_{t+1},\theta_t)
+ \hat{\beta}\,\operatorname{Var}(\theta_t)
+ \hat{\beta}^2\,\operatorname{Cov}(\theta_{t+1}\theta_t,\theta_t)
+ \operatorname{Cov}\!\Big(\frac{u_t}{p_{t-1}},\theta_t\Big).
\end{equation}

Given $u_t$ is mean zero and independent of $\theta$, the last term is zero, moreover if we assume that $\theta_t$ is iid
then every term except the variance is zero and we obtain
\begin{equation}
\hat{\beta}=\hat{\beta},
\end{equation}

which implies that every $\hat{\beta}$ leads to a self-fulfilling equilibrium.

\subsection{Multiple regressors}
We now move to the case in wich instead of a single regressor $\theta_t$ agents observe a vector of $K$ different signals $\theta_{t} = (\theta_{1,t}, \theta_{2,t}, \ldots, \theta_{K,t})$.
To be consistent with dimensions we write the PLM as
\begin{equation}
r_{t+1} =  \theta_t \beta + \varepsilon_t,
\end{equation}
where $\beta = (\beta_1, \beta_2, \ldots, \beta_K)^\top$ is a $K$-dimensional column vector of parameters.
The ALM follows accrodingly as
\begin{equation}
    r_{t+1} = \theta_{t+1} \hat{\beta}+ \theta_t \hat{\beta}+ (\theta_{t+1} \beta)(\theta_t \beta) + \frac{u_t}{p_{t-1}}.
\end{equation}

Defining 
\begin{equation}
    \Theta_t = \begin{pmatrix}
    \theta_{1} \\
     \theta_{2}\\
        \vdots \\
    \theta_{t} 
\end{pmatrix}
\quad 
R_t = \begin{pmatrix}
    r_{2} \\
     r_{3}\\
        \vdots \\
    r_{t+1}
\end{pmatrix},
\end{equation}
as the $t \times K$ matrix of regressors and the $t$-dimensional vector of returns up to time $t+1$,  the consistency condition requires

\begin{equation}
    \hat{\beta} = (\Theta_t^\top \Theta_t)^{-1} \Theta_t^\top R_{t} = \beta,
\end{equation}

and the ALM in matrix form is
\begin{equation}
    R_t = \Theta_{t+1} \hat{\beta}+ \Theta_t \hat{\beta}+ (\Theta_{t+1} \hat{\beta})(\Theta_t \hat{\beta}) + \frac{U_t}{P_{t-1}},
\end{equation}

where again we use the notation $U_t$ and $P_{t-1}$ to denote the vectors of shocks and prices up to the corresponding time periods.

Now asymptotically we can show that 

\begin{equation}
    \mathbb{E}[\Theta_t^\top R_t] = \mathbb{E}[\Theta_t^\top \Theta_{t+1}] \hat{\beta}+ \mathbb{E}[\Theta_t^\top \Theta_t] \hat{\beta} + \mathbb{E}[\Theta_t^\top (\Theta_{t+1} \hat{\beta})(\Theta_t \hat{\beta})] + \mathbb{E}\Big[\Theta_t^\top \frac{U_t}{P_{t-1}}\Big] = \mathbb{E}[\Theta_t^\top \Theta_t] \hat{\beta},
\end{equation}


which again implies that every $\hat{\beta}$ is a self-fullfilling equilibriums since
\begin{equation}
    \mathbb{E}[\hat{\beta}] = \mathbb{E}[(\Theta_t^\top \Theta_t)^{-1}] \mathbb{E}[\Theta_t^\top R_{t}] = \hat{\beta}.
\end{equation}

\section{A simple example}

Consider a reduced form model 
\begin{equation}
    r_t = \alpha r^e_t + \delta z_{t-1} + \varepsilon_t,
\end{equation}

in which all quantitites are scalars, $r_t$ is the asset return at time $t$, $r^e_t$ is the expected return formed at time $t-1$, $z_{t-1}$ is an observable signal at time $t-1$ and $\varepsilon_t \sim N(0, \sigma^2_\varepsilon)$.
The REE requires that $r^e_t = \mathbb{E}_{t-1}[r_t] = \frac{\delta}{1-\alpha}$.
We now consider the situation in which agents do not know the exact value of the parameters $\alpha$ and $\delta$ and instead form their expectations using a perceived law of motion (PLM) of the form
\begin{equation}
    r^e_t = \beta_{t-1} z_{t-1},
\end{equation}
where $\beta_{t-1}$ is the estimate of the regression coefficient obtained using data up to time $t-1$.
Substituting the PLM into the reduced form model we obtain the actual law of motion (ALM)
\begin{equation}
    r_t = \alpha \beta_{t-1} z_{t-1} + \delta z_{t-1} + \varepsilon_t.
\end{equation}
This relationships imply a map between the PLM and ALM capured by the function
\begin{equation}
    T(\beta) = \alpha \beta + \delta,
\end{equation}

and its fixed point coincides with the REE value of the regression coefficient $\beta^* = \frac{\delta}{1-\alpha}$.
We now analyze the stability of this fixed point, which depends on the mechanism used by agents to update their estimates.

\subsection{Recursive Least Squares}
The first updating mechanism we consider is Recursive Least Squares (RLS), which obtains as a recursive implementation of the standard Ordinary Least Squares (OLS).
In OLS, the agent estimates the coefficients of a linear regression model by minimizing the sum of squared residuals.
Assuming the general linear model 

\begin{equation}\label{eq:linear_model}
    y_i = \beta' x_i + \varepsilon_i,
\end{equation}

where $x_i$ is the $k \times 1$ vector of regressors, the OLS estimator can be derived from the minimization problem 

\begin{equation}
    \min_{\beta} \sum^t_{i=1} (y_i - \beta' x_i) ^2,
\end{equation}

with solution
\begin{equation}
    \hat{\beta} \equiv b = \left(\sum^t_{i=1}  x_i x_i' \right)^{-1} \sum^t_{i=1}  x_i y_i.
\end{equation}

Estimating OLS on the whole avaialble data set, is sometimes referred to as off-line estimation. 
Another approach is to estimate the model on-line, that is as soon as a nes data point arrive. 
It turns out that the estimator formula allows for a recursive formulation.
Define 
\begin{equation}
R_t = \frac{1}{t}  \left(\sum^t_{i=1}  x_i x_i' \right), \quad b_t = \frac{1}{t} R^{-1}_t \sum^t_{i=1}  x_i y_i.
\end{equation}

First we establish a recursion for $R_t$. Start from 

\begin{equation}
R_{t-1} = \frac{1}{t-1}  \left(\sum^{t-1}_{i=1}  x_i x_i' \right),
\end{equation}

then 

\begin{equation}
\frac{t-1}{t} R_{t-1} = \frac{1}{t}  \left(\sum^{t-1}_{i=1}  x_i x_i' \right)^{-1} \rightarrow \frac{t-1}{t} R_{t-1} + \frac{1}{t} x_t x_t' = R_t.
\end{equation}

That is 
\begin{equation}
    R_t = R_{t-1} + \frac{1}{t} \left(x_t x_t' - R_{t-1} \right).
\end{equation}

Similarly for $b_t$ we start from
\begin{equation}
    b_{t-1} = \frac{1}{t-1} R^{-1}_{t-1} \sum^{t-1}_{i=1}  x_i y_i,
\end{equation}

then notice that 
\begin{equation}
    (t-1) R_{t-1} = (t R_t - x_t x_t'),
\end{equation}
which implies that 
\begin{equation}
    b_{t-1}(t R_t - x_t x_t') =  \sum^{t-1}_{i=1}  x_i y_i \rightarrow b_{t-1}(1 - t^{-1} R^{-1}_t x_t x_t') = t^{-1} R^{-1}_t \sum^{t-1}_{i=1}  x_i y_i,
\end{equation}
and
\begin{equation}
    b_{t-1}(1 - t^{-1} R^{-1}_t x_t x_t') + t^{-1} R^{-1}_t x_t y_t = b_t,
\end{equation}

which can be rearranged as
\begin{equation}
    b_t = b_{t-1} + t^{-1} R^{-1}_t x_t (y_t - x_t' b_{t-1}).
\end{equation}

The appealing feature of this recursive formulation for our purpose is that it complies with the general formulation of a stochastic recursive algorithm (SRA) of the form
\begin{equation}
    \theta_t = \theta_{t-1} + \gamma_t Q(t, \theta_{t-1}, X_t),
\end{equation}
where $\theta_t$ is the parameter vector to be estimated, $\gamma_t$ is a sequence of gains, $Q(\cdot)$ is a function that depends on the time index, the current parameter estimate and the new data point $X_t$.
Under suitable assumptions about the gain sequence $\gamma_t$, and the functional form of $Q(\cdot)$, asymptotic properties are well established.
First we show how to recast the RLS updating for our simple model into the SRA form.
A small change is required to accomodate for the fact that only lagged values of the parameter estimates are used in $Q(\cdot)$.
Define $S_{t-1} = R_t$, then we obtain the updating rules
\begin{equation}
    \begin{aligned}
        \beta_t = \beta_{t-1} + t^{-1} S^{-1}_{t-1} z_{t-1} \left[(z_{t-1}(T(\beta_{t-1}) - \beta_{t-1}) + \varepsilon_t )\right], \\
        S_{t} = S_{t-1} + t^{-1}\frac{t}{t-1} (z^2_{t} - S_{t-1}),
    \end{aligned}
\end{equation}

where we used $T(\beta) = \alpha \beta + \delta$.


Now we appropriately define 
\begin{equation}
    \begin{aligned}
        \theta_t = \begin{pmatrix}
        \beta_t \\
        S_t
    \end{pmatrix} , \quad
    \gamma_t = t^{-1}, \quad
    X_t = \begin{pmatrix}
        z_t \\
        z_{t-1} \\
        \varepsilon_t
        
    \end{pmatrix}
    \end{aligned}
\end{equation}

\cite{Evans_2001} show that the stability of the SRA can be analyzed by studying the smalle associated ordinary differential equation (ODE)
\begin{equation}
    \frac{d \beta}{d \tau} = T(\beta) - \beta,
\end{equation}
where $\tau$ is notational time.
This notion, often referred to as E-stability principle, notes that the REE, which is also the fixed point of the ODE, is stable if the ODE is asymptotically stable.
In this case the stability condition is simply given by $\alpha < 1.$
In this case RLS learning converges to the REE.

\subsection{Weighted Least Squares and Constant Gain Learning}

In Weighted Least Squares (WLS), the agent estimates the coefficients of a linear regression model by minimizing a weighted sum of squared residuals.
The weights are typically chosen to reflect the relative importance or reliability of different observations. Assuming the linear model (\ref{eq:linear_model})
 the WLS estimator can be derived from the minimization problem
\begin{equation}
    \min_{\beta} \sum^t_{i=1} {w_{t,i}} (y_i - \beta' x_i) ^2,
\end{equation}

with solution
\begin{equation}
    \hat{\beta} = \left(\sum^t_{i=1} {w_{t,i}} x_i x_i' \right)^{-1} \sum^t_{i=1} {w_{t,i}} x_i y_i.
\end{equation}

In paritcular, we are interested in the case where the weights decline geometrically over time, that is
\begin{equation}
    w_{t,i} = \gamma (1 - \gamma)^{t-i}, \quad \gamma \in (0,1).
\end{equation}
Similar to the OLS case a recursive relationship can be derived.

Define 
\begin{equation}
    R_t = \gamma \sum^t_{i=1} {(1 - \gamma)^{t-i}} x_i x_i', \quad b_t = \gamma R^{-1}_t \sum^t_{i=1} {(1 - \gamma)^{t-i}} x_i y_i
\end{equation}

First we establish a recursion for $R_t$. Start from
\begin{equation}
R_{t-1} = \gamma \sum^{t-1}_{i=1} {(1 - \gamma)^{t-1-i}} x_i x_i',
\end{equation}
then
\begin{equation}
(1 - \gamma) R_{t-1} = \gamma \sum^{t-1}_{i=1} {(1 - \gamma)^{t-i}} x_i x_i' \rightarrow (1 - \gamma) R_{t-1} + \gamma x_t x_t' = R_t.
\end{equation}

That is
\begin{equation}
    R_t = R_{t-1} + \gamma \left(x_t x_t' - R_{t-1} \right).
\end{equation}

Similarly for $b_t$ we start from
\begin{equation}
    b_{t-1} = \gamma R^{-1}_{t-1} \sum^{t-1}_{i=1} {(1 - \gamma)^{t-1-i}} x_i y_i,
\end{equation}
then
\begin{equation}
    (1- \gamma) b_{t-1} = \gamma R^{-1}_{t-1} \sum^{t-1}_{i=1} {(1 - \gamma)^{t-i}} x_i y_i ,
\end{equation}

and
\begin{equation}
 b_{t-1} ( R_t - \gamma x_t x_t') =  \gamma \sum^{t-1}_{i=1} {(1 - \gamma)^{t-i}} x_i y_i \rightarrow b_{t-1} (I - \gamma R^{-1}_t x_t x_t') = \gamma R^{-1}_t \sum^{t-1}_{i=1} {(1 - \gamma)^{t-i}} x_i y_i,
\end{equation}

since \[R_{t-1} = \frac{1}{1 - \gamma} (R_t - \gamma x_t x_t').\]
Adding $\gamma R^{-1}_t x_t y_t$ to both sides we obtain
\begin{equation}
    b_{t-1} (I - \gamma R^{-1}_t x_t x_t') + \gamma R^{-1}_t x_t y_t = b_t,
\end{equation}
which can be rearranged as
\begin{equation}
    b_t = b_{t-1} + \gamma R^{-1}_t x_t (y_t - x_t' b_{t-1}).
\end{equation}

In relation to our simple model we get then the same updating rules, with the only difference that the gain is now constant and equal to $\gamma$.
Given that the gain is not decreasing and that the updating depends on the stochastic term $\varepsilon_t$, the learning rule can not converge to a fixed point.
However, under suitable conditions and for small values of $\gamma$, and under the same stability condition $\alpha < 1$ it can be shown that the learning dynamics will fluctuate around the REE and in particular that
\begin{equation}
    \beta_t \sim \mathcal{N} (\beta^*, \gamma \sigma_{\varepsilon}^2 (2 m_z (1 - \alpha))^{-1}),
\end{equation}
where $m_z = \mathbb{E}[z_t^2]$, and which makes clear that only in the limit case $\gamma \rightarrow 0$ the variance of the fluctuations vanishes and we recover convergence to the REE.

\subsection{Lasso Regression}
We now consider the case in which agents use Lasso regression to estimate the parameters of their PLM.
The LASSO (Least Absolute Shrinkage and Selection Operator) estimator is defined as the solution to a similar optimization problem

\[\hat{\beta}_{LASSO} = \arg\min_{\beta} \left\{\frac{1}{2}\left[(Y - X\beta)'(Y - X\beta)\right] + \lambda \sum_{j=1}^{k} |\beta_j|\right\},\]
where $\lambda \geq 0$ is a tuning parameter that controls the amount of regularization. The LASSO estimator can be interpreted as a trade-off between fitting the data well (minimizing the sum of squared residuals) and keeping the model simple (minimizing the absolute values of the coefficients).
In general there is no closed form solution for the LASSO estimator, since the absolute value function is not differentiable at zero. However, in the special case where $X$ is an orthonormal matrix, the LASSO estimator can be expressed in closed form.
To do so we can rewrite the optimization problem as
\[
\hat{\beta}_{LASSO} 
= \arg\min_{\beta} 
\Big\{ \frac{1}{2}\left[Y'Y - 2\beta'\underbrace{X'Y}_{\beta'_{OLS}} 
+ \beta'\underbrace{X'X}_{I}\beta\right]
+ \lambda \sum_{j=1}^{k} |\beta_j| \Big\},
\]
where the underbrackets follow from orthonormality and also notice that the term $Y'Y$ does not depend on $\beta$, so it does not change the minimizer.
Noticing that $A'A = \sum_{j=1}^{k} a_j^2$, we can rewrite the problem as
\[
\hat{\beta}_{LASSO} = \arg\min_{\beta} \Big\{ \sum_{j=1}^{k} \left( -\beta_j \beta_{OLS,j} + \frac{1}{2}\beta_j^2 + \lambda |\beta_j| \right) \Big\}.
\]

This problem is separable in the coefficients $\beta_j$, so we can solve for each coefficient independently. The optimization problem for each $j$ coefficient is
\[\min_{\beta_j} \left\{\frac{1}{2} \beta_j^2 - \beta_j \beta_{OLS,j} + \lambda |\beta_j| \right\}.\]
To solve this problem, we can consider two cases:
\begin{enumerate}
    \item $\beta_{OLS,j} > 0$, then $\beta_j$ has to be positive, otherwise we could decrease the objective by setting $\beta_j = 0$. 
    In this case, the problem becomes 
    \[\min_{\beta_j \geq 0} \left\{ \frac{1}{2} \beta_j^2 - \beta_j \beta_{OLS,j} + \lambda \beta_j \right\},\]
    and we can now differentiate and set to zero to find the minimum:
    \[\frac{\partial}{\partial \beta_j} \left( \frac{1}{2} \beta_j^2 - \beta_j \beta_{OLS,j} + \lambda \beta_j \right) = \beta_j - \beta_{OLS,j} + \lambda = 0 \Rightarrow \beta_j = \beta_{OLS,j} - \lambda.\]
    However, we also need to check that this solution is non-negative, so we have
    \[\hat{\beta}_{LASSO,j} = \max\left(0, \beta_{OLS,j} - \lambda\right).\]
    \item $\beta_{OLS,j} < 0$, then $\beta_j$ has to be negative, otherwise we could decrease the objective by setting $\beta_j = 0$. In this case, the problem becomes
    \[\min_{\beta_j < 0} \left\{ \frac{1}{2} \beta_j^2 - \beta_j \beta_{OLS,j} + \lambda (-\beta_j) \right\}.\]    
    Differentiating and setting to zero, we get
    \[\frac{\partial}{\partial \beta_j} \left( \frac{1}{2}\beta_j^2 - \beta_j \beta_{OLS,j} - \lambda \beta_j \right) = \beta_j - \beta_{OLS,j} - \lambda = 0 \Rightarrow \beta_j = \beta_{OLS,j} + \lambda.\]

    Again, we need to check that this solution is negative,
    so we have
    \[\hat{\beta}_{LASSO,j} = \min\left(0, \beta_{OLS,j} + \lambda\right).\]
\end{enumerate}

Combining both cases, we can write the LASSO estimator for each coefficient as
\begin{equation}
\hat{\beta}_{LASSO} = S(\beta_{OLS}, \lambda) = \textrm{sign}(\beta_{OLS}) \cdot \max\left(0, |\beta_{OLS}| - \lambda\right),
\end{equation}

where $S(\cdot)$ is the soft-thresholding operator.

In our simple model, we can still use the recursion for the OLS estimator, and then apply the soft-thresholding operator to obtain the LASSO estimate.
The mapping between the PLM and ALM now becomes
\begin{equation}
    \beta_{OLS} = T(\beta_{LASSO}) = \alpha \beta_{LASSO} + \delta = \alpha S(\beta_{OLS}, \lambda) + \delta,
\end{equation}
which can be solved case by case:
\begin{enumerate}
    \item If $\beta_{OLS} > \lambda > 0$, then $\beta_{LASSO} = \beta_{OLS} - \lambda \implies \beta_{OLS} = \beta^* + \frac{\alpha}{1 - \alpha} \lambda$.
    \item If $\beta_{OLS} < -\lambda < 0$, then $\beta_{LASSO} = \beta_{OLS} + \lambda \implies \beta_{OLS} = \beta^* - \frac{\alpha}{1 - \alpha} \lambda$.
    \item If $|\beta_{OLS}| \leq \lambda$, then $\beta_{LASSO} = 0 \implies \beta_{OLS} = \delta$.
\end{enumerate}

If we limit ourselves to the case $\delta > 0$, we can summarize the fixed points of the mapping as
\begin{equation}
    \beta_{OLS} = \begin{cases}
    \beta^* + \frac{\alpha}{1 - \alpha} \lambda, & \text{if } \delta > \lambda, \\
    \delta, & \text{if } \delta \leq \lambda.
    \end{cases}
\end{equation}

While this holds in equilibrium, we already discussed that with constant gain learning the estimates will fluctuate around these fixed points.
Identically to the previous case, the estimates will be distributed as a normal distribution centered around the fixed point, with variance proportional to the gain parameter $\gamma$.
So even for $\delta \leq \lambda$, there will be fluctuations in the Lasso estimate around zero, which implies that sometimes the estimates will be exactly zero (when $|\beta_{OLS}| \leq \lambda$) and sometimes they will be non-zero (when $|\beta_{OLS}| > \lambda$).
Given the distribution of $\beta_{OLS}$, we can compute the probability that the Lasso estimate is exactly zero as
\begin{equation}
    \mathbb{P}(|\beta_{OLS}| \leq \lambda) 
    = \Phi\left(\frac{\lambda - \mu_{OLS}}{\sigma_{OLS}}\right) - \Phi\left(\frac{-\lambda - \mu_{OLS}}{\sigma_{OLS}}\right),
\end{equation}
where $\mu_{OLS}$ and $\sigma_{OLS}$ are the mean and standard deviation of the distribution of $\beta_{OLS}$, and $\Phi(\cdot)$ is the cumulative distribution function of the standard normal distribution.

\bibliography{references}
\end{document}